\section{Method}
\label{sec:method}

We describe PALETTE (\S\ref{subsec:m-transform}), how it is applied during
training (\S\ref{subsec:m-training}), and the causal ablation used to
isolate texture preservation as the source of its effect
(\S\ref{subsec:m-ablation}).

\subsection{The PALETTE transform}
\label{subsec:m-transform}

Let $x$ denote the intensity volume of a training scan, min-max normalised
to $[0,1]$, and $F \subset \Omega$ its foreground voxels (an intensity
threshold; $x_v > 0.01$ in our implementation). PALETTE applies four steps
at every training iteration, illustrated end-to-end in
\cref{fig:method-pipeline} on one real CHAOS (abdominal MRI) and one real
OnHarmony (brain MRI) volume. The figure's five panels (a)--(e) are
deliberately lettered rather than numbered $1$--$5$ so they are never
confused with the four method steps below: panels are \emph{intermediate
diagnostic renderings}, one per step plus the input, not a one-to-one
visual walkthrough of the steps themselves (panel (e), the final training
input, is the only panel that shows the texture the transform actually
produces; (b)--(d) show flat preview values, defined in the caption).

\paragraph{1. Intensity partition.} We run 1-D $k$-means on the foreground
intensities $\{x_v : v \in F\}$ with $k$ sampled per iteration from
$\{2,\dots,6\}$, producing $k$ intensity classes $C_1, \dots, C_k$
(\cref{fig:method-pipeline}, panel (b)). This partition is unsupervised and
requires no anatomical label map: it groups voxels by intensity alone, so
the same operator applies unmodified to any input image regardless of its
modality or contrast --- CT, T1, T2, FLAIR, or otherwise. This step is
occasionally skipped in favour of a single global remap (step 3 below, with
$C=1$); per-step probabilities are set once and never tuned per dataset,
and are listed in \cref{sec:suppl-config}.

\paragraph{2. Spatial sub-parcellation.} Each intensity class $C_i$ is
further subdivided into spatial sub-regions by a Voronoi partition
seeded with points sampled within $C_i$ (\cref{fig:method-pipeline}, panel (c)
shows the combined effect of this step and the remap of step 3), giving
sub-regions $R_{i,1}, \dots, R_{i,m_i}$ that are contiguous in space as
well as similar in intensity. This second partition breaks the spatial
correlation of a purely intensity-based grouping: two voxels with the same
intensity in different parts of the anatomy can land in different
sub-regions and receive independent contrast perturbations. Each class is
usually, but not always, sub-split this way (\cref{sec:suppl-config}).

\paragraph{3. Signed affine remap.} For each sub-region $R$, let $\bar{x}_R$
denote its mean intensity, $\mu_R \sim U(0,1)$ a resampled target mean, and
draw $\alpha_R \sim \pm U(0.5, 2.0)$ (a magnitude in $[0.5, 2.0]$ with an
independently sampled sign). Every voxel $v \in R$ is remapped as
\begin{equation}
  y_v = \mu_R + \alpha_R \, (x_v - \bar{x}_R).
  \label{eq:affine}
\end{equation}
Because \cref{eq:affine} is a per-region affine function of the
\emph{real} intensity $x_v$, it preserves the relative structure of the
signal within each region: edges, gradients, and partial-volume transitions
between adjacent voxels are scaled, not replaced. Only the region-to-region
contrast relationship is randomized, and $\alpha_R < 0$ additionally allows
local contrast inversions that no acquisition protocol produces on its own,
widening the space of contrasts the network sees beyond what any single
real sequence could supply. Background voxels are zeroed at this step (and
remain zero through the rest of the pipeline). The result is occasionally
Gaussian-blurred, mostly not (\cref{sec:suppl-config}).

\paragraph{4. Anatomical-label affine remap (optional).} Steps 1--3 are
unsupervised, so the $k$-means/Voronoi partition need not respect a true
anatomical boundary: when two adjacent tissues share similar intensity,
step 3 can by chance assign them the same target contrast and erase that
boundary rather than randomizing it. This step corrects for that: for each
real anatomical label $\ell$ present, we resample a fresh
$(\mu_\ell, \alpha_\ell)$ and apply \cref{eq:affine} \emph{again}, restricted
to $\ell$'s voxels, \textbf{on top of the step-3 output} rather than the
original voxels (\cref{fig:method-pipeline}, panel (d)) -- so it composes
with, rather than replaces, the unsupervised remap, and $\ell$'s voxels
still inherit the internal structure steps 1--3 imposed on them. Because
label assignment is independent of the unsupervised partition, this
reintroduces a full-strength contrast discontinuity at the true label
border regardless of whether step 3 preserved it, sharpening the boundary
cue available to the network -- not the underlying anatomy. It consumes
only the labels the segmentation task already requires, so it needs no
additional annotation, and is optional: disabling it
(\texttt{label\_remap\_prob}$=0$) degrades PALETTE gracefully to steps 1--3
alone, e.g.\ on an unlabelled dataset. A final foreground z-score
normalisation follows an optional second blur pass, giving the training
input shown in panel (e) of \cref{fig:method-pipeline}. Per-step
probabilities are set once, never tuned per dataset, and listed in full in
\cref{sec:suppl-config}.

\begin{figure*}[t]
  \centering
  \footnotesize
  \setlength{\tabcolsep}{2pt}
  \newcommand{\methodrowlabel}[1]{\rotatebox{90}{\parbox{2.1cm}{\centering #1}}}
  \newcommand{\methodpanel}[1]{\includegraphics[width=3.3cm]{figures/method_panels/#1}}
  \resizebox{\linewidth}{!}{%
  \begin{tabular}{c ccccc}
    & (a) Input
    & (b) $k$-means flat remap
    & (c) + Voronoi sub-parcellation (flat remap)
    & (d) + label flat remap
    & (e) Final output (affine remap, \cref{eq:affine}) \\
    \methodrowlabel{CHAOS T2SPIR\\(abdominal MRI)}
    & \methodpanel{chaos_a} & \methodpanel{chaos_b} & \methodpanel{chaos_c}
    & \methodpanel{chaos_d} & \methodpanel{chaos_e} \\
    \methodrowlabel{OnHarmony T1w\\(brain MRI)}
    & \methodpanel{onharmony_a} & \methodpanel{onharmony_b} & \methodpanel{onharmony_c}
    & \methodpanel{onharmony_d} & \methodpanel{onharmony_e} \\
  \end{tabular}%
  }
  \caption{The PALETTE pipeline on one real CHAOS T2SPIR volume (top) and
  one real OnHarmony T1w volume (bottom); every panel is a deterministic
  function of the same real volume (no external model), min--max normalised
  to $[0,1]$ for display only. Panels (b)--(d) paint each region,
  sub-region, or label with its precomputed target mean $\mu$ from
  \cref{eq:affine}, isolating that flat value from the texture applied on
  top of it; panel (e) is the actual training input, produced by applying
  the full signed affine remap of \cref{eq:affine} with real texture (not a
  flat fill) at both the region and, where available, anatomical-label
  level.}
  \label{fig:method-pipeline}
\end{figure*}

\subsection{Training protocol}
\label{subsec:m-training}

PALETTE is applied on-the-fly as one contrast-augmentation operator inside
an existing GPU-accelerated nnU-Net augmentation
pipeline~\cite{molinier2026auglab,isensee2021nnunet}, alongside that
pipeline's standard spatial and artifact augmentations; it adds no learned
parameters and no measurable training-time overhead. A single
configuration -- the same $k$-range, sub-region count, blur schedule, and
$\alpha$ range given in \S\ref{subsec:m-transform} -- is tuned once and
reused unchanged across every dataset and modality in
\cref{sec:experiments}; no dataset receives its own hyperparameter sweep.

\subsection{Causal ablation: isolating texture preservation}
\label{subsec:m-ablation}

PALETTE differs from label-generative synthesis (e.g.\ SynthSeg) along two
axes at once: the partition it perturbs (unsupervised $k$-means $+$ Voronoi
vs.\ a supplied anatomical label map as the sole generative substrate) and
the perturbation itself (an affine remap of real intensities vs.\
resampling from a per-label Gaussian mixture). Comparing PALETTE to
SynthSeg directly therefore cannot attribute an advantage to either axis in
isolation.

We construct a controlled variant, \emph{noise-fill}, that holds the
$k$-means/Voronoi partition and the per-region statistics $(\mu_R,
\alpha_R)$ from \cref{eq:affine} exactly fixed, but fills each sub-region
with samples drawn from a Gaussian with that region's $(\mu_R, \alpha_R
\sigma_R)$ instead of remapping the real intensities inside it -- i.e.\ the
same partitioning scheme, with SynthSeg-style noise substituted for the
affine remap. Because the partition and target statistics are identical to
PALETTE's own, any difference in downstream performance between PALETTE and
noise-fill isolates the one variable that changed: whether real texture is
preserved or discarded within each region. We report this comparison
separately for structures with no tissue interface (fine internal gradients, low boundary
contrast) and structures bounded by one
(identified by a sharp, high-contrast border) in \cref{sec:experiments}, as
the texture-preservation hypothesis predicts a widening gap specifically on
the former.
